% Formula bodies can be copied individually into MathType TeX input.
\documentclass[UTF8]{ctexart}
\usepackage{amsmath,amssymb}
\begin{document}
% (1) current
\begin{equation}
I=(\mathrm{ADC}_{\mathrm{code}}-\mathrm{zero}_{\mathrm{code}})\times0.00125
\tag{1}
\end{equation}
式中，I为电机电流，单位A；ADC\_code为ADS1115单端原始码，zero\_code为零电流校准码。0.00125 A/LSB由125 μV/LSB的ADC尺度与100 mV/A的霍尔灵敏度换算得到。

% (2) rpm
\begin{equation}
n=\frac{\Delta C_{\mathrm{raw}}\times60\times1000}{4096\times100}
\tag{2}
\end{equation}
式中，n为有符号转速，单位rpm；ΔC\_raw为100 ms测速窗内的编码器计数增量；4096为每转计数，60和1000分别完成分钟及毫秒单位换算。

% (3) stats1
\begin{equation}
\mu=\frac{1}{N}\sum_{i=1}^{N}x_i,\quad\mathrm{RMS}=\sqrt{\frac{1}{N}\sum_{i=1}^{N}x_i^2},\quad\mathrm{PP}=\max(x)-\min(x)
\tag{3}
\end{equation}
式中，x\_i为窗口内第i个样本，N为样本数；μ为均值，RMS为均方根，PP为峰峰值。统计结果随输入通道采用对应的原始码或物理量单位。

% (4) stats2
\begin{equation}
\sigma^2=\frac{1}{N}\sum_{i=1}^{N}(x_i-\mu)^2,\quad\mathrm{CF}=\frac{\max|x|}{\mathrm{RMS}},\quad K=\frac{\frac{1}{N}\sum_{i=1}^{N}(x_i-\mu)^4}{\sigma^4}
\tag{4}
\end{equation}
式中，σ²为总体方差，σ为标准差，CF为峰值因子，K为Pearson峭度；max|x|为窗口内绝对峰值，μ、x\_i和N与前式相同。CF和K为无量纲量。

% (5) fft
\begin{equation}
X[k]=\frac{1}{N}\sum_{n=0}^{N-1}x[n]w[n]\exp\left(-\frac{\mathrm{j}2\pi kn}{N}\right),\quad f[k]=\frac{kf_s}{N}
\tag{5}
\end{equation}
式中，x[n]为第n个时域样本，w[n]为窗函数，N为FFT点数；j为虚数单位，k为频点索引，X[k]为经1/N缩放的复频谱，f\_s为采样率，f[k]为频率，单位Hz。

% (6) class
\begin{equation}
p(c\mid x)=\frac{\exp(z_c)}{\sum_{j=0}^{4}\exp(z_j)},\quad c^*=\mathop{\operatorname{argmax}}_c p(c\mid x)
\tag{6}
\end{equation}
式中，x为窗口特征，z\_c为类别c的原始输出，p(c|x)为softmax分数；c=0～4依次表示正常、过载、欠压、过热和电机倾斜，c*为最终选定类别。

% (7) ordinal
\begin{equation}
q_{j,k}=\operatorname{sigmoid}(s_j-\theta_{j,k}),\quad P(g=1,2,3)=(1-q_{j,1},q_{j,1}-q_{j,2},q_{j,2})
\tag{7}
\end{equation}
式中，j表示已选异常类别，s\_j为其等级潜在分数，θ\_j,k为第k个有序阈值，k=1、2；q\_j,k表示等级超过k的概率，g为严重度等级，P给出1～3级的概率分布。

% (8) grade
\begin{equation}
g_j=1+\mathbf{1}[q_{j,1}\ge0.5]+\mathbf{1}[q_{j,2}\ge0.5]
\tag{8}
\end{equation}
式中，g\_j为异常类别j的输出等级；1[·]为指示函数，条件成立取1，否则取0。两个有序阈值共同将等级分为1、2、3级。

% (9) cls_loss
\begin{equation}
L_{\mathrm{cls}}=-\frac{\sum_{i=1}^{B}w_{c_i}\log p_{i,c_i}}{\sum_{i=1}^{B}w_{c_i}}
\tag{9}
\end{equation}
式中，B为批次样本数，c\_i为样本i的真实类别，p\_i,c\_i为该类别预测分数，w\_c\_i为由训练集类别频数计算的权重；L\_cls为加权交叉熵。

% (10) ord_loss
\begin{equation}
L_{\mathrm{ord}}=-\frac{1}{2B_F}\sum_{i\in F}\sum_{k=1}^{2}\left[t_{i,k}\log q_{i,k}+(1-t_{i,k})\log(1-q_{i,k})\right]
\tag{10}
\end{equation}
式中，F为批次内具备有效等级标签的异常样本集合，B\_F为其样本数，t\_i,k为是否超过等级阈值k的二值标签，q\_i,k为真实故障类别对应的累计等级概率；L\_ord为有序分级二元交叉熵。

% (11) loss
\begin{equation}
L=L_{\mathrm{cls}}+0.5L_{\mathrm{ord}},\quad t_{i,k}=\mathbf{1}[g_i>k]
\tag{11}
\end{equation}
式中，L为故障模型总训练损失，0.5为分级损失权重，g\_i为真实故障等级。正常样本参与分类损失，异常样本在所属故障的等级头上接受分级监督；批次无有效等级标签时，分级损失记为0。

% (12) prf
\begin{equation}
P=\frac{TP}{TP+FP},\quad R=\frac{TP}{TP+FN},\quad F_1=\frac{2PR}{P+R}
\tag{12}
\end{equation}
式中，P为精确率，R为召回率，F\_1为二者的调和平均；TP、FP、FN分别为目标类别的正确识别数、误识别数和漏识别数。宏平均F\_1对五类F\_1取算术平均。

% (13) kan
\begin{equation}
y_o=\sum_i w_{b,oi}\operatorname{SiLU}(x_i)+\sum_i\sum_k w_{s,oik}B_{i,k}(x_i)
\tag{13}
\end{equation}
式中，x\_i和y\_o分别为输入与输出分量；w\_b,oi为基础分支权重，w\_s,oik为样条分支权重，B\_i,k为输入i对应的第k个B样条基函数，SiLU为基础激活函数。

% (14) pcc
\begin{equation}
\operatorname{PCC}(u,v)=\frac{(u-\bar u)\cdot(v-\bar v)}{\|u-\bar u\|\|v-\bar v\|+\varepsilon}
\tag{14}
\end{equation}
式中，u、v为参与相关计算的两个特征向量，ū、v̄为其均值，·为内积，‖·‖为欧氏范数，ε为数值稳定项；PCC为Pearson相关系数。

% (15) attention
\begin{equation}
S=\operatorname{softmax}\left(\frac{QK^T}{\sqrt{d_k}}\right),\quad G=\operatorname{softmax}(\operatorname{ReLU}(A))
\tag{15}
\end{equation}
式中，Q、K、V分别为查询、键、值矩阵，K^T为转置，d\_k为单头键向量维度；S为归一化注意力矩阵，A为相关图邻接矩阵，G为归一化图先验，ReLU用于保留非负图关系。

% (16) aggregate
\begin{equation}
H=\operatorname{softmax}[(1-\lambda)S+\lambda G]V
\tag{16}
\end{equation}
式中，λ为图先验融合系数，本模型取0.5；H为融合注意力对值矩阵V加权后的输出，softmax沿聚合节点维度归一化。

% (17) rul
\begin{equation}
\mathrm{RUL}=c_{\mathrm{terminal}}-c_{\mathrm{window,end}}
\tag{17}
\end{equation}
式中，c\_terminal为该电机最后观测老化批次编号，c\_window,end为当前历史窗口末次观测批次编号；RUL为两者之差，单位为剩余观测老化试验轮次。

% (18) mse
\begin{equation}
L_{\mathrm{MSE}}=\frac{1}{B}\sum_{i=1}^{B}(\hat y_i-y_i)^2
\tag{18}
\end{equation}
式中，B为当前批次样本数，y\_i为真实剩余轮次，ŷ\_i为预测轮次；L\_MSE为训练均方误差，单位为轮²。每轮训练损失按样本数加权汇总。

% (19) quant
\begin{equation}
q=\operatorname{clip}\left(\operatorname{round}(x/s)+z,q_{\min},q_{\max}\right),\quad x\approx s(q-z)
\tag{19}
\end{equation}
式中，x为浮点数值，q为量化整数，s为量化尺度，z为零点；round为取整，clip为范围截断，q\_min、q\_max为对应整数类型的上下界。

% (20) errors
\begin{equation}
\mathrm{MAE}=\frac{1}{n}\sum_{i=1}^{n}|\hat y_i-y_i|,\quad\mathrm{RMSE}=\sqrt{\frac{1}{n}\sum_{i=1}^{n}(\hat y_i-y_i)^2}
\tag{20}
\end{equation}
式中，n为评估窗口数，y\_i和ŷ\_i分别为窗口i的真实与预测剩余轮次；MAE为平均绝对误差，RMSE为均方根误差，二者单位均为轮。

% (21) mape
\begin{equation}
\mathrm{MAPE}=\frac{100\%}{n_+}\sum_{i\in I_+}\frac{|\hat y_i-y_i|}{|y_i|},\quad I_+=\{i:y_i\ne0\}
\tag{21}
\end{equation}
式中，I\_+为真实标签非零的评估窗口集合，n\_+为该集合样本数；MAPE为平均绝对百分比误差，单位为%。分段结果按相应寿命区间内的窗口统计。

\end{document}